\chapter{The Birthday Paradox}

\section{The Problem}
In a room with \(n\) randomly chosen people, what is the probability that
\emph{at least two} people share the same birthday?

The surprising fact (the ``birthday paradox'' \cite{wiki:birthday-problem}) is
that this probability becomes large even when \(n\) is not very big. For
example, with \(n=23\) it already exceeds \(50\%\) under a simple model.

Throughout this chapter we make the following assumptions:
\begin{itemize}
  \item there are \(365\) equally likely birthdays (ignore leap years),
  \item each person's birthday is independent of the others.
\end{itemize}

\section{How to calculate the probability}
Let \(B\) be the event that \emph{at least two} people share a birthday.
It is easier to compute the complementary event \(A\): \emph{all birthdays are
distinct}. Then
\[
  \Pr(B)=1-\Pr(A).
\]

\subsection*{Step 1: Compute \(\Pr(A)\)}
Imagine the people enter the room one by one.
\begin{itemize}
  \item The first person can have any birthday: probability \(1\).
  \item The second person must avoid the first person's birthday: probability
  \(\frac{364}{365}\).
  \item The third person must avoid the first two birthdays: probability
  \(\frac{363}{365}\).
  \item \(\cdots\)
  \item The \(n\)-th person must avoid the previous \(n-1\) birthdays:
  probability \(\frac{365-(n-1)}{365}\).
\end{itemize}

Multiplying these conditional probabilities gives
\[
  \Pr(A)
  = 1\cdot \frac{364}{365}\cdot \frac{363}{365}\cdots
  \frac{365-(n-1)}{365}
  = \prod_{k=0}^{n-1}\frac{365-k}{365}.
\]
Equivalently,
\[
  \Pr(A)=\frac{365\cdot 364\cdot 363\cdots (365-n+1)}{365^n}
  =\frac{365!}{(365-n)!\,365^n},
\]
valid for \(n\le 365\). (For \(n\ge 366\), \(\Pr(A)=0\) by the pigeonhole
principle.)

\subsection*{Step 2: Take the complement}
Therefore, the probability of at least one shared birthday is
\[
  \Pr(B)=1-\Pr(A)
  = 1-\prod_{k=0}^{n-1}\frac{365-k}{365}.
\]

\section{Approximation}
Now replace 365 by \(m\) to treat a general ``birthday space'' of size \(m\).
\[
\Pr(A) = \prod_{k=0}^{n-1}\frac{m-k}{m}.
       = \prod_{k=0}^{n-1}(1-\frac{k}{m})
       \approx \prod_{k=0}^{n-1}e^{-\frac{k}{m}}
       =e^{-\sum_{k=0}^{n-1}\frac{k}{m}}
       =e^{-\frac{n(n-1)}{2m}}
       \approx e^{-\frac{n^2}{2m}}
\]
\[
n\approx \sqrt{-\frac{2m}{\ln(p)}}
\]
For \(m=365\) and \(\Pr(A) = 0.01\), we have \(n=57.98\).
This means that, in a room with 59 people, there is only a 1\% chance that
everyone has a distinct birthday (equivalently, a 99\% chance that at least two
people share a birthday).

\subsection*{Table for \(p=0.01\) and different sample spaces \(m\)}
Using
\[
n\approx \sqrt{-\frac{2m}{\ln(p)}}
\qquad\text{with}\qquad p=0.01,
\]
we obtain the following values:
\begin{center}
\begin{tabular}{l r r}
\hline
\textbf{Case} & \textbf{\(m\)} & \textbf{\(n \approx \sqrt{-2m/\ln(0.01)}\)} \\
\hline
Birthdays (days in a year) & \(365\) & \(58\) \\
6/49 lottery & \(13{,}983{,}816\) & \(11349\) \\
EuroMillions \(5/50\) combinations & \(139{,}838{,}160\) & \(35884\) \\
\hline
\end{tabular}
\end{center}
Readers must be careful how to interpret the results before investing in lotteries.
For example, for the 6/49 lottery, after 11349 draws, there is a 1\% chance that 
all the results are distinct.

\section{Hash Collisions and the Expected Number of Collisions}
The same counting logic applies to hash functions.

Assume a hash function outputs a uniform random integer in
\(\{1,2,\dots,n\}\). Hash \(m\) independent inputs. This is the birthday problem
with ``days'' replaced by hash values.

\subsection*{Collision probability}
The probability of at least one collision among the \(m\) hash outputs is
\[
1-\prod_{k=0}^{m-1}\frac{n-k}{n},
\]
for \(m\le n\). (If \(m\ge n+1\), a collision is guaranteed by the pigeonhole
principle.)

\subsection*{Expected number of collisions}
Now we ask a different question: what is the \emph{expected number of
collisions}?

We count ``collisions'' as follows: when the \(k\)-th hash value is produced,
it is a collision if it matches \emph{any} of the previous \(k-1\) hash values.
(So if three values are equal, this counts as two collisions: the second repeat
and the third repeat.)

For the \(k\)-th output to \emph{avoid} a collision, it must avoid matching each
of the previous \(k-1\) outputs. Under the uniform model, the probability of no
match with a particular previous output is \((n-1)/n\), so the probability of
no match with all \(k-1\) previous outputs is
\[
\left(\frac{n-1}{n}\right)^{k-1}.
\]
Therefore, the probability that the \(k\)-th output \emph{is} a collision is
\[
1-\left(\frac{n-1}{n}\right)^{k-1}.
\]

Let \(X\) be the total number of collisions in this sense among \(m\) outputs.
By linearity of expectation,
\[
\mathbb{E}[X]
=\sum_{k=1}^{m}\left[1-\left(\frac{n-1}{n}\right)^{k-1}\right]
= m-\sum_{k=0}^{m-1}\left(\frac{n-1}{n}\right)^k.
\]
The remaining sum is a geometric series:
\[
\sum_{k=0}^{m-1}r^k=\frac{1-r^{m}}{1-r}\qquad (r\ne 1).
\]
With \(r=(n-1)/n\),
\[
\sum_{k=0}^{m-1}\left(\frac{n-1}{n}\right)^k
=\frac{1-\left(\frac{n-1}{n}\right)^m}{1-\frac{n-1}{n}}
=n\left(1-\left(\frac{n-1}{n}\right)^m\right).
\]
So the expected number of collisions is
\[
\boxed{
\mathbb{E}[X]
= m-n+n\left(\frac{n-1}{n}\right)^m.
}
\]
This derivation follows \cite{might:hash-collisions}.

\section{Approximation by Binomial Distribution}
There are \(n\) people and \(m\) possible birthdays.
If we pair a person with a birthday, there are \(n\times m\) possible
person--birthday pairs.

There are \(binom{n}{2}\) possible pairs of people.
So the probability that one person shares a birthday with another person is

\[p=\frac{\binom{n}{2}}{n\times m-1}
\approx \frac{\binom{n}{2}}{n\times m}
\].

Let \(X\) be the number of collisions. 

We choose \(k\) from \(n\) people. There are \(\binom{n}{k}\) different ways to
choose them. There are \(k\) ``collision'' outcomes each with probability \(p\),
and \(n-k\) ``non-collision'' outcomes each with probability \(1-p\). Therefore
the probability is
\[
\Pr(X=k)=\binom{n}{k}p^k(1-p)^{n-k}.
\]

This distribution is called the \emph{binomial distribution}.
Since there are \(n\) people, the expected number of collisions is
\(\frac{\binom{n}{2}}{m}\).

\subsection*{Poisson limit theorem (why binomial \(\to\) Poisson)}
\[
\Pr(X_n=k)=\binom{n}{k}p^k(1-p)^{n-k}.
\]
We now show that this converges to a Poisson\(\!\left(\frac{\binom{n}{2}}{m}\right)\)
probability.

Let \[
\lambda=\frac{\binom{n}{2}}{m}.
\]
\[
\binom{n}{k}
=\frac{n(n-1)\cdots(n-k+1)}{k!}
=\frac{n^k+O(n^{k-1})}{k!}.
\]
So
\[
\binom{n}{k}p^k(1-p)^{n-k}
=\left(\frac{n^k+O(n^{k-1})}{k!}\right)
\left(\frac{\lambda}{n}(1+o(1))\right)^k
(1-p)^{n-k}.
\]
The first two factors simplify to \(\frac{\lambda^k}{k!}(1+o(1))\). For the last
factor, split:
\[
(1-p)^{n-k}=(1-p)^n(1-p)^{-k}.
\]
As \(n\to\infty\), we have \((1-p)^{-k}\to 1\) and
\[
(1-p)^n
=\left(1-\frac{\lambda}{n}(1+o(1))\right)^n
\to e^{-\lambda}.
\]
Therefore,
\[
\Pr(X_n=k)=\binom{n}{k}p^k(1-p)^{n-k}\to e^{-\lambda}\frac{\lambda^k}{k!}.
\]
So \(X_n\) is well-approximated by \(\mathrm{Pois}(\lambda)\) when \(n\) is large
and \(p\) is small with \(np\approx \lambda\). This is the Poisson limit
theorem (law of rare events) \cite{wiki:poisson-limit-theorem}.

